\chapter{Background}
\label{ch:background}

% Background covers foveated vision, the saliency/scanpath lineage, foveated models and the
% \dg{} architecture. Dataset and metrics are defined where first used in \chref{methods};
% subject differences, including age, are in \chref{discussion}.

This thesis asks what happens to a model of where people look next when it is made to see the scene with
the resolution falloff of a human eye.

\section{Foveated vision}
\label{sec:foveated-vision}

Human vision is not uniform across the visual field. Acuity is highest at the \emph{fovea}, a region
roughly $2^\circ$ of visual angle in radius around the point of gaze, and falls off steeply with distance
from it into the periphery \citep{strasburger2011peripheral}. Only a small patch of any scene is seen
sharply at any instant.

This is why the eyes move. Looking at a scene is a cycle of \emph{fixations}, during which the eye is
nearly still and visual information is taken in, separated by \emph{saccades}, the fast jumps between
them, during which little is seen \citep{rayner1998eye}. Where to look next has to be chosen from what the
blurred periphery shows, together with memory of what was already inspected \citep{henderson2003gaze}.
\citet{najemnik2005} showed that human search is close to a Bayesian ideal observer that picks each
saccade to maximise the probability of finding the target, weighing every candidate fixation by what the
eccentricity-dependent visibility map would reveal from there.
\figref{one-fixation} shows one image sharp and at two foveation strengths.

\begin{figure}[tbp]
  \centering
  \includegraphics[width=\textwidth]{background/one_fixation.png}
  \caption[One image, sharp and foveated]{One image from MIT1003 \citep{judd2009mit1003}, sharp and foveated at two strengths, at the dataset's $35$ pixels per degree.
    The red cross is the fixation and the yellow box the patch magnified below. The dashed circle in
    B, of radius $3^\circ$, encloses the region the blur leaves untouched (\subsecref{sharp-radius});
    at $f_{c0} = 10$ no such region is left.}
  \label{fig:one-fixation}
\end{figure}

Detail is measured in cycles per degree, the number of repeats of a pattern that fit into one degree of
visual angle (a thumb held at arm's length covers about two degrees). The more repeats fit into that
degree, the finer the detail, and blurring an image removes the finest.

This thesis models the falloff as a cutoff spatial frequency that falls with eccentricity, following
\citet{geislerperry1998}, discarding at every point in the image the detail finer than the cutoff for
that eccentricity. The model folds into one curve the coarser retinal
sampling away from the fovea \citep{curcio1990photoreceptor} and the shrinking cortical area per degree
of visual field, called \emph{cortical magnification}
\citep{daniel1961representation, strasburger2011peripheral}. It
also treats the falloff as identical in every direction, whereas measured acuity differs between
directions in the visual field (\figref{retinal-falloff}). The shape of the falloff is set by one parameter, $e_2$, the eccentricity at which
the cutoff has halved. \chref{methods} uses $e_2 = 2.3^\circ$, the value \citet{geislerperry1998} fitted to human
contrast-sensitivity measurements.
The \citet{geislerperry1998} model is used here because one number sets its strength, the cutoff at the
point of gaze, $f_{c0}$, in cycles per degree. That number can be set to the human value, and below it to degrade
the periphery further than an eye does.
The blur models the loss of resolution only. Peripheral vision also loses in other ways, among them
crowding, the failure to identify an object when others lie close to it \citep{rosenholtz2016peripheral},
which the model leaves out (\secref{disc-limits}).

\begin{figure}[tbp]
  \centering
  \includegraphics[width=0.72\textwidth]{background/retinal_falloff.png}
  \caption[The acuity model against measured falloff]{The model of \chref{methods} against the falloff it simplifies. The coloured curves are what
    \citet{watson2014formula}'s ganglion-cell density fit implies along the four directions out from
    the fovea; the dashed one is the Geisler--Perry cutoff used here, falling more slowly because it is
    fitted to what a subject can see rather than to what the retina samples.}
  \label{fig:retinal-falloff}
\end{figure}

\section{Saliency and scanpath models}
\label{sec:models}

A \emph{saliency map} is the spatial distribution of fixations; a \emph{scanpath} is their order. Both are
targets a gaze model can be asked to predict, and \figref{saliency-vs-scanpath} shows the same data under
each reading. \dg{} shows the two are not exclusive, computing a saliency map and using it to predict the next
fixation in a scanpath, but they still set different problems.

A saliency map can be built from the data itself, by collecting every fixation many subjects made on an
image and smoothing them together (\figref{saliency-vs-scanpath}A).

Scanpaths cannot be built that way. A saliency map has one value per pixel, and every recorded fixation
is evidence about one of them. A scanpath is an ordered sequence, so the same counting would need
evidence about each sequence a subject might produce; the number of possible sequences grows with their
length, while a dataset holds a handful per image. A scanpath model
therefore has to compute the next fixation from the image and the fixations so far, rather than look it
up.

\begin{figure}[tbp]
  \centering
  \includegraphics[width=\textwidth]{background/saliency_vs_scanpath.png}
  \caption[A saliency map and one scanpath]{One MIT1003 image \citep{judd2009mit1003} under the two readings. \textbf{A.} Every fixation of
    all fifteen subjects, smoothed into a density, with order discarded. \textbf{B.} One subject's
    fixations, numbered in viewing order.}
  \label{fig:saliency-vs-scanpath}
\end{figure}

The models below use \emph{saliency map} in a second, mechanistic sense, as bottom-up conspicuity computed
from the image alone, where a region stands out because it differs from its surroundings in contrast,
colour or orientation. The brain is thought to select saccade targets from a \emph{priority}
map, that same image-driven signal plus whatever the viewer's goals make relevant and a memory of
where they have already looked \citep{fecteau2006salience, zelinsky2015priority}.
\figref{dg3-composition} shows the two maps \dg{} adds to predict one fixation. Its probability map
has the image-driven signal and a memory of the fixations already made, and
\figref{history-conditioning} shows that memory changing the prediction as the history grows. The
goals have no counterpart in it, because the subjects of MIT1003 viewed each image freely with no
task set.

\begin{figure}[tbp]
  \centering
  \includegraphics[width=\textwidth]{background/dg3_composition.png}
  \caption[\dg{} composing one prediction]{\dg{} given the first four fixations of one subject (drawn). \textbf{A.} The read-out's map,
    from the image and those fixations together. \textbf{B.} The centerbias (\secref{dg3-arch}),
    scaled by the weight the model learned. \textbf{C.} Their sum in log units, normalised. A and C
    share a colour scale; B carries its own.}
  \label{fig:dg3-composition}
\end{figure}

\dg{} reuses the convolutional read-out of the static DeepGaze line and adds conditioning on fixation
history, joining the two lines \figref{lineage} traces.

\begin{figure}[tbp]
  \centering
  \begin{tikzpicture}[
    seed/.style  = {draw, rounded corners=2pt, fill=black!8, align=center,
                    font=\footnotesize, text width=22mm, minimum height=12mm},
    sal/.style   = {draw, rounded corners=2pt, fill=blue!12, align=center,
                    font=\footnotesize, text width=22mm, minimum height=12mm},
    scan/.style  = {draw, rounded corners=2pt, fill=green!14, align=center,
                    font=\footnotesize, text width=22mm, minimum height=12mm},
    host/.style  = {draw, very thick, rounded corners=2pt, fill=blue!25, align=center,
                    font=\footnotesize\bfseries, text width=24mm, minimum height=14mm},
    flow/.style  = {-{Latex[length=2.2mm]}, thick},
    graft/.style = {-{Latex[length=2.2mm]}, densely dashed, thick, gray},
    lane/.style  = {font=\scriptsize\bfseries\itshape, anchor=center},
  ]
    % --- lane labels ---
    \node[lane] at (2.9, 5.0)  {static saliency};
    \node[lane] at (2.9, 2.0)  {scanpath generation};

    % --- root ---
    \node[seed] (r1) at (0.0, 2.5) {Feature-integration\\ theory (1980);\\ saliency map (1985)};

    % --- static-saliency lineage ---
    \node[sal]  (s1) at (2.9, 4.0)  {Itti--Koch (1998)};
    \node[sal]  (s2) at (5.8, 4.0)  {GBVS, AIM, SUN\\ (2006--09)};
    \node[sal]  (s3) at (8.7, 4.0)  {DeepGaze I--IIE\\ (2015--21)};

    % --- scanpath-generation lineage ---
    \node[scan] (p1) at (2.9, 1.0)  {winner-take-all\\ + inhibition of\\ return (1998)};
    \node[scan] (p2) at (5.8, 1.0)  {saccadic models\\ (2015)};
    \node[scan] (p3) at (8.7, 1.0)  {deep path samplers\\ (2017--20)};

    % --- confluence and what follows ---
    \node[host] (dg) at (11.9, 2.5) {DeepGaze III (2022)\\ density over the\\ next fixation};
    \node[scan] (nx) at (11.9,-1.0) {transformer \&\\ diffusion (2023--25)};

    % --- succession within the two lineages ---
    \draw[flow] (r1) -- (s1);
    \draw[flow] (r1) -- (p1);
    \draw[flow] (s1) -- (s2);   \draw[flow] (s2) -- (s3);
    \draw[flow] (p1) -- (p2);   \draw[flow] (p2) -- (p3);

    % --- where the lineages meet ---
    \draw[flow]  (s3) -- (dg);
    \draw[flow]  (p3) |- (nx.west);
    \draw[graft] (p3) -- (dg);
    \draw[graft] (dg) -- (nx);
  \end{tikzpicture}
  \caption[Two lines of work, and where they meet]{Static saliency learns where people look on average; scanpath generation produces the order.
    \dg{} inherits the frozen-backbone read-out of the first and the history conditioning of the
    second. Solid arrows mark descent, dashed ones an idea carried across lines.}
  \label{fig:lineage}
\end{figure}

\subsection{From hand-crafted saliency to learned features}
\label{subsec:saliency-lineage}

\citet{treisman1980feature} set out the picture later models inherit. Simple
visual features are registered across the whole field at once, and attention is then
directed to one location at a time. Guided Search added that the second stage is steered
by the first rather than deployed blindly \citep{wolfe1989guided}.
\citet{koch1985shifts} turned Treisman's account into a mechanism whose parts are still in use: the
feature maps are combined into one topographic \emph{saliency map}, a winner-take-all
network selects its peak, and inhibiting that location briefly --- what was later named
\emph{inhibition of return} --- lets the next-most-active one win. They describe shifts of
attention rather than eye movements, but this is how later models read an ordered
sequence of fixations off a single map, as \citet{itti1998saliency} first implemented.

These are \emph{hand-crafted} models, in which a person chooses which features to compute and how
to weight them, and nothing is fitted to eye-tracking data. Because the
result emits both a map and a path, it is at once a saliency model and a saccade
generator. Later hand-crafted models kept the architecture and varied how activation
on the map was computed \citep{harel2006gbvs, bruce2009saliency, zhang2008sun}.

Deep Gaze~I \citep{kummerer2015deepgaze1} took the ImageNet feature maps of
\citet{krizhevsky2012alexnet}, left them untouched, and fitted only a small read-out on top
of them to human fixations. It reported the best results on the MIT Saliency Benchmark at the
time, without a face or text detector anywhere in the model, because the semantic content that draws
fixations was already in the features, which hand-crafted models had to supply through
purpose-built detectors. That fixed the recipe
the whole line then followed, and the one used here, freezing a pretrained network's
features and training only a small read-out on top of them
\citep{kummerer2016deepgaze2, linardos2021deepgaze2e}. Every model in this line,
hand-crafted or learned, emits a single static map, predicting \emph{where} fixations
land but not their order.

\subsection{Scanpath-prediction models}
\label{subsec:scanpath-lineage}

Scanpath models predict the fixation sequence itself. The earliest is the winner-take-all
with inhibition of return of \subsecref{saliency-lineage}, but it always jumps to the
current peak, so it produces one fixed path per image and reproduces neither the
variability of human gaze nor its \emph{oculomotor} regularities (how far and in which
direction people tend to move their eyes, independently of what the image contains). Those
regularities predict fixation locations at least as well as image salience does
\citep{tatler2009prominence}. Le~Meur and Liu built them in directly, drawing candidates from the saliency map
\emph{reweighted} by the measured human saccade amplitude and direction distributions and by a memory
term that suppresses locations just visited, rather than always taking the peak
\citep{lemeur2015saccadic}.

Deep models learn the dynamics instead, from saliency volumes sampled over time
\citep{assens2017saltinet} to \emph{inverse reinforcement learning} for goal-directed search
\citep{yang2020predicting}, which infers from human search paths what reward those paths were
maximising and then acts on it. These differ in machinery but all emit paths rather than a probability
for an observed fixation, so they are scored by how closely a sampled path resembles a human
one, through order-alignment metrics such as ScanMatch \citep{cristino2010scanmatch} and Sequence Score
\citep{yang2020predicting}.
\citet{kummerer2021state} show that path-to-path similarity can rank a \emph{wrong} model above the
true generating distribution. Human scanpaths on one image disagree with each other
(\figref{scanpath-variability}), so any single human path
is one draw from a broad distribution, and so is any single model path. Two draws from
the same distribution routinely look unalike, while a model that collapses onto one
stereotyped path can resemble each individual
human more closely than another human does. Scoring the density instead asks how much probability the
model placed where people looked, which variability cannot distort in that direction.

\begin{figure}[tbp]
  \centering
  \includegraphics[width=\textwidth]{background/scanpath_variability.png}
  \caption[Scanpath variability between subjects]{The first five of the fifteen subjects on two stimuli, each panel one subject's fixations in
    viewing order with the first in white. \textbf{Top:} the same object, in different orders and
    with different fixation counts. \textbf{Bottom:} a dispersed stimulus, where the fixations land
    on different parts of the image from subject to subject.}
  \label{fig:scanpath-variability}
\end{figure}

\dg{} \citep{kummerer2022deepgaze3} closes that gap. Instead of producing a path, it produces a
probability for every location the next fixation could land on, given the image and the fixations so
far. That probability can be read at the place a person looked, so the model is scored on what it
predicted there rather than on how closely one of its paths resembles one of theirs. It is the model this thesis instruments, and \chref{methods} describes how it
is built. \figref{history-conditioning} shows what conditioning on the history does to the prediction.

\begin{figure}[tbp]
  \centering
  \includegraphics[width=\textwidth]{demos/priority_evolution/priority_evolution_stim0091.png}
  \caption[Conditioning on fixation history]{The probability map \dg{} predicts after $1$, $3$ and $4$ fixations of one subject's scanpath,
    with that scanpath overlaid. The path is the subject's, not a model sample, so each panel is where
    the model expected that subject to go next, and the gap between the two is what the likelihood of
    \subsecref{metrics} measures.}
  \label{fig:history-conditioning}
\end{figure}

\subsection{Current state (2023--2025)}
\label{subsec:contemporary}

Transformers \citep{vaswani2017attention, dosovitskiy2021vit} and diffusion models
\citep{ho2020ddpm} became standard tools across computer vision, and scanpath researchers
took them up as they arrived.

Gazeformer predicts
goal-directed search attention with a transformer that encodes the search target
through a language model, which lets it generalise to targets never seen in training
\citep{mondal2023gazeformer}. HAT unifies top-down search and bottom-up free-viewing
in a single transformer, predicting a dense heat-map per fixation where earlier models chose
from a coarse grid of cells; it is also the only model in this section with an explicit foveal
mechanism, which \secref{foveated-models} takes up \citep{yang2024hat}. EyeFormer uses a transformer policy trained by
reinforcement learning and adds personalisation, adapting to a handful of scanpaths from one subject
\citep{jiang2024eyeformer}. TPP-Gaze models fixation position and \emph{duration} jointly with neural
temporal point processes \citep{damelio2025tppgaze}, so a fixation's timing is a distribution. The
model here leaves timing out. ScanDiff uses a
diffusion model to generate \emph{diverse} scanpaths, taking the spread of plausible
paths as the thing to predict \citep{cartella2025scandiff}.

The transformer backbones carry a data cost. A transformer has none of the locality and translation
equivariance a convolution builds in, so it has to learn them, and \citet{dosovitskiy2021vit} report that
their model generalises worse than a comparable convolutional network when trained on ImageNet alone
and overtakes it only after pretraining on a dataset an order of magnitude larger. Several of the models
above are built for visual search rather than free viewing, a difference \secref{position} returns to.

\section{Foveated models}
\label{sec:foveated-models}

Most deep saliency and scanpath models read the whole image at uniform resolution. The few that
reintroduce a fovea differ in \emph{where} they attach it, and the subsections below take those places
in pipeline order.

\subsection{At the input}
\label{subsec:fov-input}

Input-space models degrade the image before the network sees it, by resampling or by blurring.

\emph{Resampling} takes closely
spaced samples near the fovea and progressively sparser ones outward, as the retina's own receptors are
spaced. Laying those samples back out on an even grid magnifies the centre and squeezes the periphery, so
straight lines bend and a pixel no longer sits where it came from. \citet{dacosta2024retinal} apply this
to a recognition network rather than a gaze model, but the transform is the same one a foveated gaze
model would inherit.

A space-variant \emph{blur} leaves every pixel exactly where it was and changes only
how sharp the picture is there, as SemBA-FAST does for semantic search \citep{luzio2025semba}. The search
model of \citet{yang2020predicting} uses a two-level version of the same idea, a sharp patch at the
fixation laid over a blurred copy of the whole frame. This thesis is input-space and uses the blur.

STAR-FC \citep{wloka2018starfc} is the closest prior work, the one input-space model that foveates for
free viewing, the task studied here. It re-applies a retinal transform at each new fixation and takes the
next fixation as the maximum of a priority map built from conspicuity and a history of visited locations.
It selects that maximum rather than emitting a density, so it emits ordered sequences and cannot be
scored by likelihood (\tabref{model-capabilities}). Its transform is the same space-variant blur, applied
anisotropically where the falloff here is identical in every direction. STAR-FC then splits the blurred
frame, running a low-level saliency algorithm over the periphery and a deep network over the centre; the
model here reads the whole foveated frame with one network.

\subsection{At the features}
\label{subsec:fov-features}

Feature-space models act on the network's internal representations. HAT reads
them at two resolutions, coarse over the whole frame and fine at each fixation, in what its authors call
a \enquote{simplified foveated retina} \citep{yang2024hat}. Because that stage sits \emph{after} the backbone,
the backbone still receives every detail, and the constraint governs only how the resulting features are
pooled.

\subsection{At the decision}
\label{subsec:fov-decision}

The SceneWalk lineage \citep{engbert2015scenewalk, schwetlick2019ccn, schwetlick2023diss} constrains the
selection step itself. SceneWalk drives its attention map by saliency
multiplied by a Gaussian window centred on the current gaze (\enquote{spatially-limited (foveated)
access}, so detail near fixation contributes more than the periphery), and subtracts a second map that
builds up at visited
locations to produce inhibition of return. Because the model emits a density, its parameters can be
fitted and its variants compared by likelihood \citep{schuett2017likelihood}, the same footing \dg{} is
evaluated on here.

The window weights the saliency signal inside a cognitive dynamical model, so the image is never
degraded and only the weight of each location in the decision changes.

\tabref{model-capabilities} puts the models of this chapter side by side.

\begin{table}[tbp]
  \centering
  \caption[What the models of this chapter predict, and where they foveate]{The models of this chapter,
    in order of publication, each name carrying the citation it is discussed under.
    \emph{Predicts} separates a
    single static map per image from an ordered sequence. \emph{Likelihood} marks the models that assign
    a probability to each observed fixation; the rest emit sampled paths and are compared by how
    closely a path resembles a human one (\subsecref{scanpath-lineage}). \emph{Foveation} names the
    stage at which a model imposes an eccentricity-dependent
    resolution limit, and \emph{follows gaze} marks the ones that re-centre that limit on the current
    fixation.}
  \label{tab:model-capabilities}
  \footnotesize
  \setlength{\tabcolsep}{2.5pt}
  \begin{tabular}{lccccc}
    \toprule
    model & task & predicts & likelihood & foveation & follows gaze \\
    \midrule
    Itti--Koch \citeyearpar{itti1998saliency}     & free viewing         & where + order & ---        & ---      & ---        \\
    Le~Meur \& Liu \citeyearpar{lemeur2015saccadic} & free viewing       & where + order & ---        & ---      & ---        \\
    SceneWalk \citeyearpar{engbert2015scenewalk, schwetlick2019ccn} & free viewing & where + order & \checkmark & decision & \checkmark \\
    DeepGaze I--IIE \citeyearpar{kummerer2015deepgaze1, linardos2021deepgaze2e} & free viewing & where & \checkmark & ---      & ---        \\
    SaltiNet \citeyearpar{assens2017saltinet}     & free viewing         & where + order & ---        & ---      & ---        \\
    STAR-FC \citeyearpar{wloka2018starfc}         & free viewing         & where + order & ---        & input    & \checkmark \\
    Yang et al.\ \citeyearpar{yang2020predicting} & search               & where + order & ---        & input    & \checkmark \\
    \dg{} \citeyearpar{kummerer2022deepgaze3}     & free viewing         & where + order & \checkmark & ---      & ---        \\
    Gazeformer \citeyearpar{mondal2023gazeformer} & search               & where + order & ---        & ---      & ---        \\
    HAT \citeyearpar{yang2024hat}                 & free viewing, search & where + order & \checkmark & features & \checkmark \\
    EyeFormer \citeyearpar{jiang2024eyeformer}    & free viewing         & where + order & ---        & ---      & ---        \\
    SemBA-FAST \citeyearpar{luzio2025semba}       & search               & where + order & \checkmark & input    & \checkmark \\
    TPP-Gaze \citeyearpar{damelio2025tppgaze}     & free viewing         & where + order & \checkmark & ---      & ---        \\
    ScanDiff \citeyearpar{cartella2025scandiff}   & free viewing, search & where + order & ---        & ---      & ---        \\
    \midrule
    this thesis                  & free viewing         & where + order & \checkmark & input    & \checkmark \\
    \bottomrule
  \end{tabular}
\end{table}

\section{Position of this thesis}
\label{sec:position}

\dg{}'s own authors raised the absence of foveation. Introducing the model,
\citet{kummerer2019ccn} write that \enquote{one might argue that our model is limited by the
fact that it is not foveated}. They answered it in the same paper and again in the journal version, at
the read-out, where several spatial priority channels do not perform substantially better than one
\citep{kummerer2019ccn, kummerer2022deepgaze3}. The input
was sharp in every variant, so what a foveated input does was left open, and it is the gap this thesis
occupies.

The successors of \subsecref{contemporary} change the generator (transformer,
diffusion, point process), and what they add over \dg{} is reach: a fixation history
without a fixed window, image regions that interact without depth, and in TPP-Gaze's case
fixation durations. The mechanism of \chref{methods} attaches ahead of the backbone, so it is
compatible with any of them.

Free viewing removes the goal input to target selection (\secref{models}), though not the subject's
contribution altogether, since faces, text and people are fixated more than image contrast predicts
\citep{judd2009mit1003}. Constraining what the stimulus delivers therefore bears directly on this task,
in a way it plausibly would not on search, where a model could still reason about where the target ought
to be.

The last row of \tabref{model-capabilities} pairs input foveation that follows gaze with free viewing
and a likelihood score, which no row above combines.

That fixes what the experiment asks of the model. \dg{} is trained and evaluated on input no human
subject could receive, with every pixel sharp, including the ones far enough from the modelled gaze that a
person would see only a blur. Restricting the input to what a foveated subject could resolve asks
whether that detail carries any of the model's fit to human fixations, and degrading the periphery
further asks how much the model draws on the coarse peripheral signal a viewer chooses saccades from
\citep{najemnik2005}; if it does, the cost should grow as the falloff steepens. \chref{methods} builds
the arms, and \chref{results} reports what the data show.
